\begin{abstract}
\noindent
Digital signatures provide authenticity, but they concentrate trust: a single private key
held by a single party is a single point of failure. Threshold signature schemes
distribute this trust across $n$ parties so that any $k$ of them can jointly produce a
valid signature, while fewer than $k$ gain nothing useful even if they cooperate. This
report surveys the topic of threshold signatures through three research papers that
together cover the foundational constructions, their security proofs, and practical
extensions.

The first paper we study is Shoup's ``Practical Threshold Signatures'' (Eurocrypt 2000),
which presents the first RSA-based threshold signature scheme that is simultaneously
non-interactive, has compact shares, and comes with a rigorous security proof in the
random oracle model. The second paper, Boldyreva's ``Threshold Signatures, Multisignatures and Blind
Signatures Based on the Gap-Diffie-Hellman-Group Signature Scheme'' (PKC 2003),
constructs a threshold scheme over GDH groups that avoids both the trusted dealer and
the zero-knowledge proofs required by Shoup, at the cost of requiring pairing-friendly
groups and a stricter $t < n/2$ corruption bound.
The third paper, Komlo and Goldberg's ``FROST: Flexible Round-Optimized Schnorr Threshold
Signatures'' (2020), addresses round efficiency by reducing threshold Schnorr signing to
two rounds (or one with preprocessing), standardised as RFC 9591 and deployed in
real-world systems such as Zcash.

Together, these papers trace a line from basic constructions to more practical and
expressive schemes, and show how the standard toolkit of Shamir secret sharing, Lagrange
interpolation, and zero-knowledge proofs can be assembled to achieve a property that
initially seems hard: splitting signing authority without any single party ever holding
the full key.
\end{abstract}
